\section*{Chapter \ref{ch:iv2:commodities-and-energy} --- Commodities and Energy}
\markboth{\protect\omTocMark{}{Answers: Commodities and Energy}}{\protect\omTocMark{}{Answers: Commodities and Energy}}

\begin{solution}{iq:iv2:commodities-and-energy:1}
It is contango, and you lose \$4 a barrel.

The future has to converge to spot by expiry. If spot stays at 70, the future falls from 74 to 70.

\iqwrong{Saying ``spot did not move, so I made nothing.'' You did not hold spot. Your return is the change in the future, and the slope of the curve is part of it.}
\end{solution}

\begin{solution}{iq:iv2:commodities-and-energy:2}
The upper limit is 88. At 92, buy spot, store it and sell the future, locking in \$4 a unit.

The full cost of carry is $80 \times 1.05 + 4 = 88$. Borrow 80 and buy; in a year, deliver against the future at 92, pay 84 on the loan and 4 for storage.

\iqwrong{Saying the fair value is 88. It is a ceiling. With inventory scarce, the future can sit far below it.}
\end{solution}

\begin{solution}{iq:iv2:commodities-and-energy:3}
It is \$25 a barrel. Refiners are naturally long it.

Three barrels of crude cost 240. They yield two barrels of gasoline at $2.40 \times 42 = 100.80$ each, and one of heating oil at $2.70 \times 42 = 113.40$. So the products bring 315, a margin of 75 on three barrels.

\iqwrong{Forgetting the units. Products quote per gallon, crude per barrel. Subtract before converting, and you get nonsense.}
\end{solution}

\begin{solution}{iq:iv2:commodities-and-energy:4}
You make nothing on the parallel move, and \$1,000 on the second. You are betting on tighter supply, meaning a deeper backwardation.

The position only sees the spread, which goes from 2 to 3. That is \$1 on 1,000 barrels.

\iqwrong{Thinking this position is long oil. It has no outright exposure to a parallel move, though front contracts tend to move more than back ones, so the spread is not fully neutral in practice.}
\end{solution}

\begin{solution}{iq:iv2:commodities-and-energy:5}
The spark spread is \$25 per MWh, which is \$40,000 for the block, and the clean spark spread is \$13.

Fuel is $7 \times 5 = 35$, so $60 - 35 = 25$. Power is $100 \times 16 = 1{,}600$ MWh. Carbon is $0.4 \times 30 = 12$.

\iqwrong{Treating that as profit. It leaves out start-up cost, maintenance and the fact that the plant may not run for every hour.}
\end{solution}

\begin{solution}{iq:iv2:commodities-and-energy:6}
It is about 11.9\% a year.

With $F = S e^{(r + u - y)T}$, we have $y = r + u - \ln(F/S) = 0.06 - \ln(66/70) = 0.06 + 0.0588 = 0.1188$.

The convenience yield is the value of having the metal rather than a promise of it: you can keep a plant running and meet a sudden order. It is high when inventories are low, as a shortage makes the last ton worth a lot. It is close to zero when the warehouses are full.

\iqwrong{Reading the backwardation as a forecast that the price will fall by 4. A curve this shape is a price for holding physical stock now. It does not by itself say where the spot will go.}
\end{solution}

\begin{solution}{iq:iv2:commodities-and-energy:7}
The one-month contango cannot go above about 1.17, given access to storage. Backwardation has no such limit.

Full carry is $50e^{0.04/12} + 1 = 51.17$. Above that, you buy, store and sell forward at a profit.

Backwardation would need the reverse: sell the metal, buy the future, and take delivery later. Only someone who owns inventory can do that, and a holder who needs the stuff to run a plant will not. So nothing forces it, and the curve is capped on the contango side and open on the other.

There are two exceptions. Limits to the top vanish if storage is full, as nobody can do the carry trade. And a non-storable, such as power, has no carry at all, so neither bound applies.

\iqwrong{Assuming the two sides are symmetric because the formula looks symmetric. A short physical position cannot be built from nothing, so cost-of-carry is an inequality.}
\end{solution}

\begin{solution}{iq:iv2:commodities-and-energy:8}
Commodity A loses about 11.3\% in a year. Gas earns exactly zero, however steep its curve looks month to month.

\textbf{Commodity A.} Each month the contract you hold starts 1\% above spot and falls to it. So the monthly return is $1/1.01 - 1$, and the year gives $1 - 1.01^{-12} = 11.3\%$ of loss. It is a real return because the expected spot does not rise.

\textbf{Gas.} The forward for each month sits at that month's price and expires there. You earn nothing on any contract, as it never moves. The curve's slope is entirely the calendar, and buying the next contract in a rising part is not a cost: it will be worth the same price when it expires.

Over a year, the front-to-next slopes also sum to zero, since the seasons repeat.

So a naive roll yield, front over second, would call gas a good carry in winter and a poor one in summer. That is wrong, and to measure carry you compare the same calendar month across years.

\iqwrong{Sorting commodities by front-to-second slope and calling the top the high-carry names. In seasonal markets you are sorting by the date on the calendar.}
\end{solution}

\begin{solution}{iq:iv2:commodities-and-energy:9}
The hedge ratio is 0.686, or about 163 contracts, and it removes 64\% of the variance.

The minimum-variance ratio is $h = \rho\sigma_s/\sigma_f = 0.8 \times 3/3.5 = 0.686$. On 10 million gallons, that is 6.857 million gallons of heating oil, or $6{,}857{,}143/42{,}000 = 163$ contracts.

The variance left is $1 - \rho^2 = 36\%$. So the standard deviation drops only by 40\%. A correlation of 0.8 sounds high, but it leaves 60\% of the risk.

\iqwrong{Using one contract per 42,000 gallons, a ratio of 1. That takes on more heating oil risk than the airline has jet fuel risk, and the leftover variance would be larger than at the minimum.}
\iqfollow{Why would the airline hedge with heating oil at all?}{It trades with far more liquidity than a jet fuel contract does. The price is the basis risk it keeps.}
\end{solution}

\begin{solution}{iq:iv2:commodities-and-energy:10}
The one-month future has 50.8\% volatility today, and the one-year future has 8.1\%. The average for the one-year option is 29.7\%, and for the one-month option 55.3\%.

Spot volatility is $0.6$ at $\tau = 0$. So $0.6\,e^{-2/12} = 0.508$ and $0.6\,e^{-2} = 0.081$.

An option on a future to expiry $T$ takes the average variance as $\tau$ runs from $T$ to 0:
\[
\bar\sigma^2 = \frac{0.36}{T}\int_0^T e^{-4\tau}\,\mathrm{d}\tau = 0.36\,\frac{1 - e^{-4T}}{4T}.
\]
At $T = 1$, $\bar\sigma = 0.6 \times 0.495 = 0.297$. At $T = 1/12$, $\bar\sigma = 0.553$.

This is the \index{Samuelson effect}Samuelson effect. A shock to spot gets largely undone by supply responses in time, so a far contract hardly reacts. As it gets close to expiry, it starts to feel each shock. That makes the mean-reversion rate and volatility level visible in the curve of implied volatilities.

\iqwrong{Quoting today's volatility of the one-year future, 8.1\%, for the one-year option. Over the year that contract ages into the front, so that number is far too low.}
\end{solution}

\begin{solution}{iq:iv2:commodities-and-energy:11}
The averaged option is worth about 1.86, against about 3.22 for the European. So it costs about 58\%.

Averaging over the whole life cuts the effective variance to a third. So the effective volatility is $40\%/\sqrt{3} = 23.1\%$ on the same time. With the at-the-money rule, $0.4 \times 70 \times 0.231 \times \sqrt{1/12} = 1.87$, and Black's formula gives 1.86. The plain option has $0.4 \times 70 \times 0.40 \times \sqrt{1/12} = 3.23$, and Black's gives 3.22.

Commodity consumers buy and sell the physical all through a month. So an average price is what they actually pay. And it makes the settlement hard to manipulate, as nobody can push one print to swing the outcome.

\iqwrong{Using $\sigma/2$ or $\sigma/\sqrt{2}$. A running average's variance is $\sigma^2T/3$, and that comes from integrating $(T - t)^2/T^2$.}
\iqfollow{What if the average is over 22 daily fixings?}{The variance factor is $(n + 1)(2n + 1)/(6n^2) = 0.356$, slightly above $1/3$, so the volatility is about 23.9\%.}
\end{solution}

\begin{solution}{iq:iv2:commodities-and-energy:12}
It is just open: about 20 cents a barrel, before quality, port and other fees.

Sale minus purchase minus freight: $77 - 73 - 3.50 = 0.50$. Financing the cargo costs $73 \times 0.05 \times 30/365 = 0.30$, which leaves 0.20.

To lock it in, buy the cargo, and sell the European benchmark for the delivery date. The US price risk is gone as soon as you buy. Freight is fixed by booking the ship or by a freight contract.

\textbf{What remains.}
\begin{itemize}
\item Quality: the crude grade may not fetch the benchmark price.
\item The benchmark's differential against the actual sale may change.
\item Timing: delay, demurrage, and weather. The hedge's date is fixed, and the ship's arrival is not.
\item Execution: if many ships head for the same port, the local price at delivery falls. The arbitrage closes on its own.
\end{itemize}

\iqwrong{Treating a futures-to-futures spread as if it were a profit. The two contracts are for different places. The difference means nothing until you've paid to move the barrel.}
\end{solution}

\begin{solution}{iq:iv2:commodities-and-energy:13}
The prices are \$0, \$10, \$40 and \$90 at 20, 60, 100 and 115 MW. At 140 demand cannot be met, so the price goes to whatever cap applies, or some load is cut. With 60 MW of wind, the prices fall to \$0, \$0, \$10 and \$40.

Stack the plants by marginal cost, and the marginal plant sets the price. Capacity up to the demand point: 30, 70, 110, 130. So 20 MW is within wind, 60 within nuclear, 100 within gas and 115 within peakers.

With 30 more MW at zero cost, all the stack shifts right: 60, 100, 140, 160. So 100 falls to nuclear and 115 to gas. That is the \index{merit-order effect}merit-order effect. More zero-cost supply lowers prices in general, including prices earned by the wind itself.

\textbf{Negative prices.} Power can't be stored at scale, and supply must equal demand each moment. Some plants pay to stay on, because shutting down and restarting is costly. And subsidised wind and solar may earn more than the negative price, so they keep producing. So they will bid below zero to keep running, and the price clears there.

\begin{dated}{2026-11}{Negative hours in Germany}\label{dat:iv2:commodities-and-energy:neg}
The German regulator counted negative wholesale prices in 573 hours of 2025 and 457 hours of 2024. Check the latest count before you cite one.
\end{dated}

\iqwrong{Averaging plant costs. A market with marginal pricing pays every generator the price of the marginal unit, and that's how the wind and nuclear plants cover their fixed costs.}
\iqfollow{A baseload forward is 50, and peak is 70, with 80 peak hours and 88 off-peak hours a week. What is off-peak?}{$(168 \times 50 - 80 \times 70)/88 = 31.82$.}
\end{solution}

\begin{solution}{iq:iv2:commodities-and-energy:14}
The option is worth about 6.94 per MWh, of which 5 is intrinsic and 1.94 is time value. At a correlation of 0.9, it falls to about 5.56. Hedge by selling about 0.74 MWh of power and buying about 4.52 MMBtu of gas.

Let the strike leg be $F_2 = 7 \times 5 = 35$. Then \index{Margrabe formula}Margrabe uses the volatility of the ratio,
\[
\sigma = \sqrt{0.6^2 + 0.4^2 - 2 \times 0.5 \times 0.6 \times 0.4} = 0.529,
\]
with $d_1 = (\ln(40/35) + \sigma^2 T/2)/(\sigma\sqrt{T}) = 0.637$ and $d_2 = 0.372$. So $V = 40\,\Phi(d_1) - 35\,\Phi(d_2) = 6.94$.

At a correlation of 0.9, the ratio's volatility is 0.297, so the value is mostly intrinsic. And as correlation goes to one with equal volatilities, it is exactly intrinsic.

\textbf{Hedge.} The deltas are $\Phi(d_1) = 0.738$ on power and $\Phi(d_2) = 0.645$ on the gas leg. That means $0.645 \times 7 = 4.52$ MMBtu of gas per MWh. As with any option you are long, you re-hedge by selling high and buying low, which earns the realised spread volatility.

\iqwrong{Using power's volatility alone. It's the volatility of the spread that matters, so high correlation makes the plant worth less.}
\iqfollow{What is it worth if power falls to 30?}{About 1.48: a plant out of the money today is still worth something because it can wait, which is why one does not retire it at the first bad month.}
\end{solution}

\begin{solution}{iq:iv2:commodities-and-energy:15}
Nothing stops it. With storage full, the carry bound fails, and a holder who cannot take delivery must pay somebody to take the contract. A price of $-38$ with the next month at 20 means storage for the month is worth about \$58 a barrel.

If you can store, $F_1 \ge F_2 - c$ holds: at a price below that, you buy, store a month and sell the next contract. So the lower limit on the front contract is the next contract less the cost of carry. At a storage cost $c$ of \$58, $20 - 58 = -38$.

Normally, $c$ is a small number. When capacity is full, it is not a price you can pay at all, so the bound does not exist. So the front is free to go wherever it clears.

Three sets of holders are forced to sell. A long-only fund, including retail products, that has to exit by its roll date. A speculator with no way to take delivery. And anyone whose risk limits require closing the position.

The buyers who can stand the delivery are few and hold storage. They name their price.

\begin{dated}{2026-11}{The day this happened}\label{dat:iv2:commodities-and-energy:neg-wti}
On 20 April 2020 the expiring May WTI contract settled at $-\$37.63$, with the next contract at 20.43, a calendar spread of about $-\$58$. The next-widest front-to-second gap in the public daily series from 1985 to 2024 is \$8.49.
\end{dated}

\iqwrong{Saying a price cannot fall below zero because oil has value. The future is a claim to take delivery at one place and date. When you can't take it, the obligation has a negative value.}
\iqfollow{Why did it not hit other contracts as hard?}{They had time. The second month had a whole month in which the storage could clear, and the producers could cut output.}
\end{solution}

\begin{solution}{iq:iv2:commodities-and-energy:16}
The intrinsic value is 0.72: inject in month 3 and withdraw in month 5. With the reveal in month 4, the tank is worth 0.87, so the option to act later adds 0.15.

\textbf{Intrinsic.} Find the best pair of dates, buy at the earlier and sell at the later. The pair is month 3 at 2.85 and month 5 at 3.60, and the costs are 0.03, so $3.60 - 2.85 - 0.03 = 0.72$. A single cycle is all that fits: the tank holds one unit, and no other pair adds value on top. Enumeration confirms it.

\textbf{Extrinsic.} You inject in month 3. In month 4, you see month 5's price, and enumerating every action in every state confirms the plan below.
\begin{itemize}
\item At 4.10, wait: you get $4.10 - 2.85 - 0.03 = 1.22$.
\item At 3.10, withdraw at month 4's price of 3.40: you get $3.40 - 2.85 - 0.03 = 0.52$.
\end{itemize}
So the value is $(1.22 + 0.52)/2 = 0.87$. That is a floor at 3.40 on the sale, and the forward at 3.60 is unchanged in expectation.

In practice, you lock in intrinsic with a static hedge, and then roll as the curve moves, each time switching into a better spread. Every switch locks in more, and that is \index{rolling intrinsic}rolling intrinsic. It doesn't count all the value, but it gives a lower bound that you can bank.

\iqwrong{Reporting the intrinsic alone. A tank is a set of spread options on the forward curve, and much of a real one's value can be in the flexibility. If nothing moves, the extra value is zero, and if volatility is high, it can be the main part.}
\iqfollow{What limits a real tank?}{The rate of injection depends on fill level, and the gas needs a minimum working level. Those limits turn this into a dynamic program with fill as a state.}
\end{solution}

\begin{solution}{iq:iv2:commodities-and-energy:17}
The margin call is \$200 million, and the present value of the gain on your sales is about \$157 million. So the one-for-one hedge is too big, even in the best case. A tailed hedge would be about 79\%.

\textbf{Flat move.} Each year's barrels gain \$20 million, which you get only when you deliver them. At 5\%, with the cash arriving mid-year, the total is
\[
\sum_{k = 0}^{9} 20 \times e^{-0.05(k + 0.5)} = \$157.4 \text{ million},
\]
against \$200 million paid now. So you are \$42.6 million worse off on a mark-to-market basis, and the cash call comes at once while the gains stretch over a decade.

\textbf{Tail the hedge.} Scaling the stack to $157.4/200 = 78.7\%$ makes it neutral to this move. The cash problem is then smaller, but it stays: you still pay \$157 million today against gains that come later. The firm has to have that credit line, signed in advance.

\textbf{Long end moves 30\%.} Long-dated prices usually move less than the front. If the average sale gains only \$6 instead of \$20, you gain about $0.3 \times 157 = \$47$ million on the sales and still lose \$200 million on the stack. A hedge ratio from a flat move is then wildly too large, so the right ratio should come from how each part of the curve actually moves.

\textbf{Contango.} At 50 cents a month on 10 million barrels, you lose \$5 million in every roll, which is about \$60 million a year at the start, falling as the stack shrinks.

So the same position can have a sound economic value and still die of funding. A lender who sees margin calls in the first year does not see the gains in the tenth.

\iqwrong{Saying ``Long oil against a short sale of oil, so it is hedged.'' Risk comes from three things: the curve can change shape, the cash comes at different times, and the rolling cost depends on curve shape.}
\end{solution}

\begin{solution}{iq:iv2:commodities-and-energy:18}
At flat volatility, the call strike is about 84.30. With skew it falls to about 75.43.

\textbf{Flat.} Under Black with zero rates, the put at 60 and 30\% costs 3.74. Solve for the call strike with the same premium: 84.30. So the producer's revenue is floored at 60 and capped at 84.30.

\textbf{Skew.} The put at 36\% costs 5.11, and a call at 26\% must collect 5.11, so the strike is 75.43. The structure protects the same floor, and gives up nearly nine dollars of upside, because the producer sells cheap vol and buys dear.

The producer gives up the upside above the call strike. ``Costless'' means no cash up front. The cost is in that upside, plus the dealer's margin in the strikes.

At the limit, put and call at the same strike $K = F$ give a synthetic short forward with zero premium.

\iqwrong{Treating zero premium as zero cost. Any fair-priced collar has zero expected profit before the dealer's margin. And with skew against you, the terms get worse without any premium appearing in the deal.}
\iqfollow{What does the bank do on its side?}{It holds the short put and long call, so it is long the future's delta. It sells futures against that and manages the skew exposure it picked up from many such deals.}
\end{solution}

\begin{solution}{iq:iv2:commodities-and-energy:19}
Kirk gives about 7.46, and a normal spread model gives about 7.70. A simulation gives 7.47. So Kirk is the one to trust here, and the simulation is how you check.

\textbf{Kirk.} Treat $F_2 + K$ as one lognormal asset with volatility scaled by $w = F_2/(F_2 + K) = 80/95 = 0.842$. Then
\[
\sigma = \sqrt{0.35^2 + (0.842 \times 0.3)^2 - 2 \times 0.9 \times 0.35 \times 0.842 \times 0.3} = 0.165.
\]
And $V = F_1\Phi(d_1) - (F_2 + K)\Phi(d_2)$, with $d_1, d_2$ from $\ln(100/95)$. This gives 7.46. With $K = 0$, it is exact: that is Margrabe.

\textbf{Normal spread.} The dollar volatilities are $35$ and $24$, so the spread's is
\[
\sqrt{35^2 + 24^2 - 2 \times 0.9 \times 35 \times 24} = 17
\]
a year, and $12.02$ over six months. With forward spread $x = 5$,
\[
V = 12.02\,\bigl(\varphi(5/12.02) + (5/12.02)\Phi(5/12.02)\bigr) = 7.70.
\]
It is about 3\% off here, as the true spread between two lognormals is not normal.

\textbf{Check.} Simulate both futures with the given correlation, and take the mean of the payoff. Kirk loses accuracy with high strikes relative to $F_2$, or very high correlation.

\iqwrong{Using a single spread volatility from the first model across strikes. The correlation that matches market spread options changes with strike, so the correlation skew is a real risk in these books, and low-volatility spreads hide it.}
\end{solution}

\begin{solution}{iq:iv2:commodities-and-energy:20}
The intrinsic value is \$70,400. As one monthly decision it is worth about \$108,000, and as 22 daily decisions about \$196,000. I would pay something well below the daily figure, about \$130,000 to \$160,000 at most, for reasons below.

\textbf{Intrinsic.} The strike leg is $8 \times 5 = 40$, so the spread is \$2. There are $100 \times 16 \times 22 = 35{,}200$ MWh. So $2 \times 35{,}200 = \$70{,}400$.

\textbf{Monthly.} Margrabe on 42 against 40, with volatility of the ratio
\[
\sqrt{0.25 + 0.1225 - 2 \times 0.6 \times 0.5 \times 0.35} = 0.403
\]
and $T = 1/12$, gives \$3.07 per MWh, so \$108,000.

\textbf{Daily.} Each day, the ratio's volatility is $\sqrt{0.81 + 0.25 - 2 \times 0.5 \times 0.9 \times 0.5} = 0.781$, at $T = 1/8$. That gives \$5.57 per MWh, so \$196,000.

The daily value is larger for two reasons. You can skip the days when the spread is negative, so you hold many options and not one. And day-ahead prices are more volatile than monthly forwards.

\textbf{Hedges.} Against the monthly value, sell $0.683 \times 35{,}200 = 24{,}059$ MWh of power and buy $0.641 \times 8 \times 35{,}200 = 180{,}501$ MMBtu of gas. Against the daily value, the deltas are lower: 21,948 MWh and 145,138 MMBtu. Both are less than a full position since the plant does not run in every state.

\textbf{The fee.} The daily number assumes you can run or not without cost, and capture each day's price. Against that:
\begin{itemize}
\item Each start costs fuel and wear, and minimum run times stop you from taking only the good hours.
\item The day-ahead volatility is not monetised by delta-hedging monthly forwards. You hold the difference between monthly and daily exposure as risk, and hedging it needs liquid daily products.
\item Outages can leave you short power on the worst day, when your hedge most needs the plant.
\item The hourly or daily correlation estimate is very uncertain, and 0.5 against 0.6 moves the value a lot.
\end{itemize}
So I would pay above the monthly number only for what I could hedge or have reason to believe, and leave a clear margin below the top of the range. The numbers here are illustrations of how to think, not rules.

\textbf{Forward at \$32.} Then the spread is $-8$. The intrinsic value is zero. The monthly value collapses to about \$0.04 per MWh, so about \$1,500. The daily value is about \$1.16, so about \$41,000. Far out of the money, nearly all the value is in day-to-day volatility. So a plant that looks dead on the monthly curve still has real value for the spikes.

\iqwrong{Paying the option value computed from one forward volatility for all 22 days. This counts neither the difference between monthly and daily decisions, nor the operating limits.}
\iqfollow{Who is your natural counterparty on this?}{A plant owner who wants a fixed income and no market risk. You are buying its optionality, and it is selling you its volatility for cash.}
\end{solution}
